% GENERATED FILE. Edit canonical lessons, then run npm run generate:books.
% canonical-source-hash: fnv1a64:6890826c3a1b06fa
% canonical-lessons: ES-C676-turismo, ES-C676-mirador, ES-C676-turistico, ES-C676-folleto, ES-C676-regatear

\chapter{Turismo}
\label{ch:es-turismo}
\hlchaptermodality{676}

\begin{chapteropening}
\textbf{By the end of this chapter:} \emph{I can talk about tourism, pick up a leaflet, find a viewpoint and haggle at a market.}

Complete the last lesson of chapter 676: I can talk about tourism, pick up a leaflet, find a viewpoint and haggle at a market.
\end{chapteropening}

\section[el turismo]{el turismo — tourism; a private car (Spain, on forms)}

\label{lesson:ES-C676-turismo}



\begin{quote}
Before the new one: say the Spanish for the arcade, then the Spanish for the paving stone.
\end{quote}

\subsection*{You'll want to know: el turismo}
\textbf{el turismo} — \textquotedblleft{}tourism; a private car (Spain, on forms)\textquotedblright{}.

\begin{quote}
El turismo es importante aquí.
\end{quote}

\begin{quote}
\emph{Tourism is important here.}
\end{quote}

\emph{Hacer turismo} — to go sightseeing.

\begin{cousinweb}
\emph{turismo} was borrowed from English \emph{tourism} (EtymDB), built on \emph{tour}, from French \emph{tour} \textquotedblleft{}a turn, a round trip\textquotedblright{}.
\end{cousinweb}

\begin{grammarlens}[title={What you've built}]
One more word for telling, explaining and describing at B1.
\end{grammarlens}

\subsection*{Your turn}
\begin{itemize}
\raggedright
  \item \emph{Say it:} \emph{el turismo}
  \item \emph{Say it:} \emph{el turismo}, once more
  \item \emph{Say it:} a sentence of your own with \emph{el turismo}
  \item \emph{Recall:} say the Spanish for the arcade, then the Spanish for the paving stone, then say \emph{el turismo} again
\end{itemize}

\begin{tcolorbox}[breakable,colback=teal!4,colframe=teal!35!black,title={Before you move on}]
What does el turismo mean? (\textquotedblleft{}Tourism; a private car (Spain, on forms)\textquotedblright{}.) Say it once more.
\end{tcolorbox}
\section[el mirador]{el mirador — the viewpoint; the glazed balcony}

\label{lesson:ES-C676-mirador}



\begin{quote}
Before the new one: say the Spanish for the paving stone, then the Spanish for tourism.
\end{quote}

\subsection*{You'll want to know: el mirador}
\textbf{el mirador} — \textquotedblleft{}the viewpoint; the glazed balcony\textquotedblright{}.

\begin{quote}
Desde el mirador se ve toda la ciudad.
\end{quote}

\begin{quote}
\emph{The whole city can be seen from the viewpoint.}
\end{quote}

Built on \emph{mirar}: a place for looking.

\begin{cousinweb}
\emph{mirador} is built on \emph{mirar} \textquotedblleft{}to look\textquotedblright{}, which is inherited from Latin \emph{mīrārī} \textquotedblleft{}to wonder at\textquotedblright{}.
\end{cousinweb}

\begin{grammarlens}[title={What you've built}]
One more word for telling, explaining and describing at B1.
\end{grammarlens}

\subsection*{Your turn}
\begin{itemize}
\raggedright
  \item \emph{Say it:} \emph{el mirador}
  \item \emph{Say it:} \emph{el mirador}, once more
  \item \emph{Say it:} a sentence of your own with \emph{el mirador}
  \item \emph{Recall:} say the Spanish for the paving stone, then the Spanish for tourism, then say \emph{el mirador} again
\end{itemize}

\begin{tcolorbox}[breakable,colback=teal!4,colframe=teal!35!black,title={Before you move on}]
What does el mirador mean? (\textquotedblleft{}The viewpoint; the glazed balcony\textquotedblright{}.) Say it once more.
\end{tcolorbox}
\section[turístico]{turístico — tourist (as an adjective); touristy}

\label{lesson:ES-C676-turistico}



\begin{quote}
Before the new one: say the Spanish for tourism, then the Spanish for the viewpoint.
\end{quote}

\subsection*{You'll want to know: turístico}
\textbf{turístico} — \textquotedblleft{}tourist (as an adjective); touristy\textquotedblright{}.

\begin{quote}
Es una zona muy turística.
\end{quote}

\begin{quote}
\emph{It is a very touristy area.}
\end{quote}

\emph{Clase turista} — economy class.

\begin{cousinweb}
\emph{turístico} is built on \emph{turista} \textquotedblleft{}tourist\textquotedblright{}, which comes from English \emph{tourist}, from French \emph{tour} \textquotedblleft{}a turn\textquotedblright{}.
\end{cousinweb}

\begin{grammarlens}[title={What you've built}]
One more word for telling, explaining and describing at B1.
\end{grammarlens}

\subsection*{Your turn}
\begin{itemize}
\raggedright
  \item \emph{Say it:} \emph{turístico}
  \item \emph{Say it:} \emph{turístico}, once more
  \item \emph{Say it:} a sentence of your own with \emph{turístico}
  \item \emph{Recall:} say the Spanish for tourism, then the Spanish for the viewpoint, then say \emph{turístico} again
\end{itemize}

\begin{tcolorbox}[breakable,colback=teal!4,colframe=teal!35!black,title={Before you move on}]
What does turístico mean? (\textquotedblleft{}Tourist (as an adjective); touristy\textquotedblright{}.) Say it once more.
\end{tcolorbox}
\section[el folleto]{el folleto — the leaflet, the brochure}

\label{lesson:ES-C676-folleto}



\begin{quote}
Before the new one: say the Spanish for the viewpoint, then the Spanish for tourist.
\end{quote}

\subsection*{You'll want to know: el folleto}
\textbf{el folleto} — \textquotedblleft{}the leaflet, the brochure\textquotedblright{}.

\begin{quote}
Coge un folleto en la oficina de turismo.
\end{quote}

\begin{quote}
\emph{Take a leaflet at the tourist office.}
\end{quote}

Hotels and museums give them out for free.

\begin{cousinweb}
\emph{folleto} came from Italian \emph{foglietto} \textquotedblleft{}a little sheet\textquotedblright{} (EtymDB), built on \emph{foglio}, from Latin \emph{folium} \textquotedblleft{}leaf\textquotedblright{}.
\end{cousinweb}

\begin{grammarlens}[title={What you've built}]
One more word for telling, explaining and describing at B1.
\end{grammarlens}

\subsection*{Your turn}
\begin{itemize}
\raggedright
  \item \emph{Say it:} \emph{el folleto}
  \item \emph{Say it:} \emph{el folleto}, once more
  \item \emph{Say it:} a sentence of your own with \emph{el folleto}
  \item \emph{Recall:} say the Spanish for the viewpoint, then the Spanish for tourist, then say \emph{el folleto} again
\end{itemize}

\begin{tcolorbox}[breakable,colback=teal!4,colframe=teal!35!black,title={Before you move on}]
What does el folleto mean? (\textquotedblleft{}The leaflet, the brochure\textquotedblright{}.) Say it once more.
\end{tcolorbox}
\section[regatear]{regatear — to haggle, to bargain}

\label{lesson:ES-C676-regatear}



\begin{quote}
Before the new one: say the Spanish for tourist, then the Spanish for the leaflet.
\end{quote}

\subsection*{You'll want to know: regatear}
\textbf{regatear} — \textquotedblleft{}to haggle, to bargain\textquotedblright{}.

\begin{quote}
En el mercado puedes regatear un poco.
\end{quote}

\begin{quote}
\emph{At the market you can haggle a little.}
\end{quote}

Shops have fixed prices; at a market stall you can try.

\begin{cousinweb}
\emph{regatear} is usually traced, through older \emph{recatear} (EtymDB), to Vulgar Latin \emph{recaptāre}, built on Latin \emph{captāre} \textquotedblleft{}to try to seize\textquotedblright{}.
\end{cousinweb}

\begin{grammarlens}[title={What you've built}]
One more word for telling, explaining and describing at B1.
\end{grammarlens}

\subsection*{Your turn}
\begin{itemize}
\raggedright
  \item \emph{Say it:} \emph{regatear}
  \item \emph{Say it:} \emph{regatear}, once more
  \item \emph{Say it:} a sentence of your own with \emph{regatear}
  \item \emph{Recall:} say the Spanish for tourist, then the Spanish for the leaflet, then say \emph{regatear} again
\end{itemize}

\begin{tcolorbox}[breakable,colback=teal!4,colframe=teal!35!black,title={Before you move on}]
What does regatear mean? (\textquotedblleft{}To haggle, to bargain\textquotedblright{}.) Say it once more.
\end{tcolorbox}
